%%%PAGE 103 rot=0 legibility=good summary=实验：用双缝干涉测光的波长（原理、Δx公式、数据处理、误差分析、注意事项）
%%%BLOCK id=103-01 ch=15 sec=用双缝干涉测光的波长 kind=method exp=1 alt=none cont=none y=0.08-0.38
\textbf{用双缝干涉测光的波长.}

\textbf{实验原理与操作}
\[ \Delta x=\frac{L\lambda}{d} \]
\begin{itemize}
\item $L$为双缝与屏的距离$L$
\item $d$为$S_1,S_2$间距$d$
\item $\Delta x$：相邻两条明条纹间距.
\item $\lambda$：入射光波长
\end{itemize}
\[ \Delta x=\frac{x_2-x_1}{n-1} \]

\textbf{数据处理}

$\lambda=\dfrac{d}{L}\Delta x$\quad 计算波长，重复测量，计算，求出波长的平均值.
%%%END
%%%BLOCK id=103-02 ch=15 sec=用双缝干涉测光的波长 kind=method exp=1 alt=none cont=none y=0.40-0.78
\textbf{误差分析}
\begin{itemize}
\item 双缝到光屏的距离$L$的测量存在误差.
\item 干涉条纹未调到最清晰的程度
\item 分划板中心刻线与干涉条纹不平行，中心刻线没有恰好位于亮条纹中心.
\item 测量多条亮条纹间距间的距离时读数不准确，此问题中条纹数未数清. % CHECK: 原文"间距间的距离"，疑为"间的距离"
\end{itemize}

\textbf{注意事项.}
\begin{itemize}
\item 单缝和双缝要平行，中心位于遮光筒的轴线上.
\item 调节测量头时，应使分划板中心刻线和亮条纹的中心对齐，记下此时手轮上的读数，转动手轮，使分划板中心刻线和另一亮条纹的中心对齐，记下此时手轮上的读数，两次读数之差就表示这两条亮条纹间的距离.
\item 不能直接测$\Delta x$. 要测多条亮条纹间距再计算得到$\Delta x$,
\item 白光的干涉观察到的是彩色条纹，其中白色在中央，红色在最外层
\end{itemize}
%%%END
%%%PAGEEND 103
